\documentclass[10pt]{amsart}
\usepackage[a4paper,margin=22mm]{geometry}
\usepackage[no-math]{fontspec}
\setmainfont{Latin Modern Roman}[SmallCapsFont={lmromancaps10-regular.otf}]
\usepackage{amsmath,amssymb,amsthm,mathtools,centernot,tensor,hyperref,graphicx,etoolbox}
\usepackage[nointegrals]{wasysym}
\newcommand{\No}[1]{}
\emergencystretch=3em
\allowdisplaybreaks
\newtheorem{theo}{Theorem}[section]
\newtheorem{lemm}[theo]{Lemma}
\newtheorem{coro}[theo]{Corollary}
\newtheorem{prop}[theo]{Proposition}
\theoremstyle{definition}
\newtheorem{rema}[theo]{Remark}
\numberwithin{equation}{section}
\newcommand{\oast}{\circledast}
%\newcommand{\ocircle}{\medcirc} utilisation nécessaire du package wasysym... laisser le temps de trouver une solution pour ce package.

\newcommand{\noopsort}[1]{}


\newcommand{\BC}{\mathbb{C}}
\newcommand{\BD}{\mathbb{D}}
\newcommand{\BN}{\mathbb{N}}
\newcommand{\BR}{\mathbb{R}}
\newcommand{\BS}{\mathbb{S}}
\newcommand{\BT}{\mathbb{T}}
\newcommand{\vF}{{\mathbf{F}}}
\newcommand{\vG}{{\mathbf{G}}}
\newcommand{\vV}{{\mathbf{V}}}
\newcommand{\vX}{{\mathbf{X}}}
\newcommand{\CB}{\mathcal{B}}
\newcommand{\CD}{\mathcal{D}}
\newcommand{\CF}{\mathcal{F}}
\newcommand{\CH}{\mathcal{H}}
\newcommand{\CJ}{\mathcal{J}}
\newcommand{\CL}{\mathcal{L}}
\newcommand{\CM}{\mathcal{M}}
\newcommand{\CN}{\mathcal{N}}
\newcommand{\CR}{\mathcal{R}}
\newcommand{\CS}{\mathcal{S}}
\newcommand{\CU}{\mathcal{U}}
\newcommand{\CW}{\mathcal{W}}

\newcommand{\FS}{\mathfrak{S}}

\newcommand{\al}{\alpha}
\newcommand{\ga}{\gamma}
\newcommand{\de}{\delta}
\newcommand{\ep}{\varepsilon}
\newcommand{\ph}{\varphi}
\newcommand{\ta}{\tau}
\renewcommand{\th}{\theta}
\newcommand{\si}{\sigma}
\newcommand{\rh}{\rho}
\newcommand{\ch}{\chi}

\newcommand{\De}{\Delta}
\newcommand{\Ph}{\Phi}

\newcommand{\pl}{\partial}
\newcommand{\wt}{\widetilde}
\newcommand{\wh}{\widehat}

%\newcommand{\Ck}[1]{\left\{#1\right\}}
\newcommand{\Dk}[1]{\left[#1\right]}
\newcommand{\K}[1]{\left(#1\right)}
\newcommand{\Br}[1]{\langle #1 \rangle}
%\newcommand{\Abs}[1]{\left|#1\right|}
\newcommand{\abs}[1]{\left|#1\right|}
\renewcommand{\No}[1]{\left\| #1 \right\|}
\newcommand*\bra[1]{\langle#1\rangle}


\newcommand{\I}{\infty}
%\newcommand{\eqv}{\Leftrightarrow}
%\newcommand{\ol}{\overline}
\DeclareMathOperator{\Tr}{Tr}

\newcommand{\sd}{\langle \nabla \rangle}
\newcommand{\sdi}{\sd^{-1}}
\newcommand{\bx}{\langle x \rangle}
\newcommand{\lxr}{\langle \xi \rangle}

\newcommand{\tx}{{t,x}}
\newcommand{\Sa}{{\FS^\al}}
\newcommand{\Sxtx}{{\FS^{2\al'}_{x \to (t,x)}}}
\newcommand{\Stxx}{{\FS^{2\al'}_{(t,x) \to x}}}
\newcommand{\st}{\star}
\newcommand{\ft}{{\frac{1}{4}}}
\newcommand{\fh}{\frac{3}{4}}

\newcommand{\R}{\mathbb{R}}
\newcommand{\tw}{\frac{1}{2}}
\newcommand{\na}{\nabla}
\newcommand{\oc}{\ocircle}
\newcommand{\os}{\oast}
\newcommand{\ls}{\lesssim}


%%%%%%%%%%%%%%%%%%%%%%%%%%%%%%%%%%%%%%%%%%%%%%%%%%%%%%%%%%%%%%%%%%%%%%%%%%%%%%%%%%%%%%%%%%%%%%%%%%%%%%%%%%%
\graphicspath{{./figures/}}

\newcommand*{\mk}{\mkern -1mu}
\newcommand*{\Mk}{\mkern -2mu}
\newcommand*{\mK}{\mkern 1mu}
\newcommand*{\MK}{\mkern 2mu}

\hypersetup{urlcolor=purple, linkcolor=blue, citecolor=red}


\newcommand*{\relabel}{\renewcommand{\labelenumi}{(\theenumi)}}
\newcommand*{\romanenumi}{\renewcommand*{\theenumi}{\roman{enumi}}\relabel}
\newcommand*{\Romanenumi}{\renewcommand*{\theenumi}{\Roman{enumi}}\relabel}
\newcommand*{\alphenumi}{\renewcommand*{\theenumi}{\alph{enumi}}\relabel}
\newcommand*{\Alphenumi}{\renewcommand*{\theenumi}{\Alph{enumi}}\relabel}


\let\oldtilde\tilde
\renewcommand*{\tilde}[1]{\mathchoice{\widetilde{#1}}{\widetilde{#1}}{\oldtilde{#1}}{\oldtilde{#1}}}
\let\oldhat\hat
\renewcommand*{\hat}[1]{\mathchoice{\widehat{#1}}{\widehat{#1}}{\oldhat{#1}}{\oldhat{#1}}}
\title[Integer-order multilinear density estimates]{The integer-order s=1 multilinear density Strichartz estimate for iterated Duhamel operators in Schatten--Sobolev classes}
\author{Sonae Hadama}
\date{Annales Henri Lebesgue8(2025),181–218; DOI10.5802/ahl.233}
\hypersetup{hidelinks}
\AtBeginDocument{\renewcommand{\qedhere}{}}
\begin{document}
\maketitle
In this paper, we study the following Hartree equation in three-dimensional space:
\begin{align}\label{NLH}
\tag{NLH} 
i \pl_t \ga &= [- \Delta + w \ast \rho_\ga, \ga], \quad \ga:\BR \to \CB(L^2_x).
\end{align}
This is a nonlinear evolution equation for operator-valued functions. We denote the set of all bounded linear operators on $L^2_x := L^2(\BR^3)$ by $\CB(L^2_x)$, convolution in space by $\ast$ and commutator by $[\cdot, \cdot]$. We assume that $w$ is a given finite signed Borel measure on $\BR^3$. For $A \in \CB(L^2_x)$, $\rh_A(x) : \BR^3_x \to \BC$ is the density function of $A$, that is, $\rh_A(x):=k(x,x)$, where $k(x,y)$ is the integral kernel of $A$. When $w = \pm \de$, we should call~\eqref{NLH} the cubic nonlinear Schr\"odinger equation; however, we consistently call~\eqref{NLH} the Hartree equation in the rest of this paper.
We denote the Schatten $\al$-class on a Hilbert space $H$ by $\FS^\al(H)$, that is,
\begin{align}
\|A\|_{\FS^\al}:= \K{\Tr(|A|^\al) }^{\frac{1}{\al}},
\end{align}
where $\Tr$ is the trace in $H$. We write $\FS^\al:=\FS^\al(L^2_x)$. We refer to~\cite{2005Simon} for more details on the Schatten classes. Note that $\Tr A$ is the total number of particles contained in the system for any density operator $A$, and $\Tr A = \I$ at least formally if $A \in \FS^\al \setminus \FS^1$ for some $\al >1$.
We write $U(t):=e^{it\De}$. Define $A_\st B := A B A^*$ for two operators $A$ and $B$. For $s \ge 0$ and $\al \in [1,\I]$, we define the Schatten--Sobolev space $\CH^{s,\al}$ by
\begin{align}
\|A\|_{\CH^{s,\al}} := \left\|\sd^s A \sd^s\right\|_\Sa.
\end{align}
\setcounter{section}{2}
\section{Preliminaries}\label{sect3}

In this section, we consider general $d$-dimensional spaces because there is no benefit to limiting our argument to three-dimensional space.

\subsection{Christ--Kiselev type lemma} 

First, we give a Christ--Kiselev type lemma in Schatten classes; however, the following result is already obtained in~\cite{1970Go-Kre} before~\cite{2001CK}. Therefore, we should call it Gohberg--Kre\u{\i}n theorem. We obtain the following result by applying~\cite[Theorem~7.2]{2003BS} (but originally by~\cite[Theorem~III.6.2]{1970Go-Kre}), setting $E_t = F_t = P_{L^2((a,t);\,L^2_x)}$ in their notations, where $P_A$ is a projection to $A$. Note that the integral kernel of $\BD = \CJ^{E,F}_\th \BT$ is $\th(t-\ta)K(t,\ta)$, where $\th(x)=\ch_{\{x \,\ge\,0\}}$.

\begin{theo}[Gohberg--Kre\u{\i}n] \label{SCK}
Let $d \ge 1$. Let $-\I \leq a < b \leq \I$, $1 < p < \I$ and $\al \in (1,\I)$. Let $(a,b)^2 \ni (t,\ta) \mapsto K(t,\ta)\in \CB(L^2_x)$ be strongly continuous. Assume that we can write $\BT \in \CB(L^2_\tx)$ by
\begin{align}
(\BT g)(t,x) = \int_a^b K(t,\ta)g(\ta)d\ta.
\end{align}
If $\BT \in \FS^\al(L^2_\tx)$, then $\BD$ defined by
\begin{align}
(\BD g)(t,x) = \int_a^t K(t,\ta)g(\ta)d\ta
\end{align}
is also in $\FS^\al(L^2_\tx)$, and
\begin{align}
\|\BD\|_{\FS^\al\left(L^2_\tx\right)} \le C_p \|\BT\|_{\FS^\al\left(L^2_\tx\right)}
\end{align}
holds.
\end{theo}



\subsection{Duality principle and its application} 
In this subsection, we assume all Hilbert spaces are separable and infinite-dimen\-sional. Let $H, K$ be complex Hilbert spaces. Let $\al \in [1,\I]$ and $A: H \to K$ be compact. We define Schatten $\al$-norm by
\begin{align}
\|A\|_{\FS^\al(H\to K)} :=
\begin{cases}
&\K{\Tr_H\left(|A^*A|^{\frac{\al}{2}}\right)}^{\frac{1}{\al}} \text{ if } 1 \le \al <\I,\\
&\|A\|_{\CB(H \to K)} \text{ if } \al = \I.
\end{cases}
\end{align}
Note that $\FS^\al(H \to H)=\FS^\al(H)$. H\"older's inequality for Schatten norm is well-known. If $\al, \al_0, \al_1 \in [0,\I]$ satisfy $\frac{1}{\al}=\frac{1}{\al_0}+\frac{1}{\al_1}$, then it holds that
\begin{align}\label{Holder}
\|A_0 A_1\|_{\FS^{\al}(H)} \le \|A_0\|_{\FS^{\al_0}(H)} \|A_1\|_{\FS^{\al_1}(H)}.
\end{align}
Note that~\eqref{Holder} implies the following. If $\al, \al_0, \al_1 \in [0,\I]$ satisfy $\frac{1}{\al}=\frac{1}{\al_0}+\frac{1}{\al_1}$, then it follows that
\begin{align}\label{g Holder}
\|A_0 A_1\|_{\FS^{\al}(H_0\to H_2)} \le \|A_0\|_{\FS^{\al_0}(H_1 \to H_2)} \|A_1\|_{\FS^{\al_1}(H_0 \to H_1)}.
\end{align}
In fact, by the singular value decomposition, we can write $A_0 = U \wt{A_0}$, where $\|\wt{A_0}\|_{\FS^{\al_0}(H_1)} = \|A_0\|_{\FS^{\al_0}(H_1 \to H_2)}$ and $U:H_1 \to H_2$ is unitary. In the same way, $A_1 = \wt{A_1} V$, where $\|\wt{A_1}\|_{\FS^{\al_1}(H_1)} = \|A_1\|_{\FS^{\al_1}(H_0 \to H_1)}$ and $V:H_0 \to H_1$ is unitary. Therefore, we obtain
%transformation de plusieurs en une
\begin{multline}
\|A_0 A_1\|_{\FS^\al(H_0 \to H_2)}\\
\begin{aligned}
 &= \|U \wt{A_0} \wt{A_1} V\|_{\FS^\al(H_0 \to H_2)} = \|\wt{A_0} \wt{A_1}\|_{\FS^\al(H_1)} \\
&\le \|\wt{A_0}\|_{\FS^{\al_0}(H_1)} \|\wt{A_1}\|_{\FS^{\al_1}(H_1)} = \|A\|_{\FS^{\al_0}(H_1 \to H_2)} \|A_1\|_{\FS^{\al_1}(H_0 \to H_1)}.
\end{aligned}
\end{multline}

For any $\al \in [1, \I]$, define
\begin{align}
\FS^\alpha &:= \FS^\alpha\left(L^2_x \to L^2_x\right),
\quad \FS^\al_{t,x} := \FS^\al\left(L^2_{t,x} \to L^2_\tx\right), \\
\FS^\al_{x \to (t,x)} &:= \FS^\al\left(L^2_x \to L^2_{t,x}\right),
\quad \FS^\alpha_{(t,x) \to x} := \FS^\al\left(L^2_{t,x} \to L^2_x\right).
\end{align}
Let $I \subset \R$ be an interval. Let $A(t) \in \CB(L^2_x)$ for all $t \in I$ and $\sup_{t \,\in\,I} \|A(t)\|_{\CB} < \I$. Assume that $I \ni t \mapsto A(t) \in \CB(L^2_x)$ be strongly continuous. We define $A^\ocircle : L^2_x \to C_t(I,L^2_x)$ and $A^\oast : L^1_t(I,L^2_x) \to L^2_x$ by
\begin{align}
\left(A^\ocircle u_0\right)(t,x) &:= (A(t)u_0)(x), \\
\left(A^\oast f\right)(x) &:= \int_I A(\ta)^* f(\ta,x) d\ta.
\end{align}
Note that $A^\ocircle$ and $A^\oast$ are formally adjoint to each other. The following lemma is very important for our analysis, but it is essentially the same thing as~\cite[Lemma~3]{2017FS}.
\begin{lemm}[Duality principle] \label{DP}
Let $p,q,\al \in [1,\I]$. The following are equivalent:
\begin{enumerate}
\item\label{lemm3.2.1} For any $\ga \in \FS^\al$,
\begin{align}\label{SS}
\left\|\rh(A(t)_\st \ga)\right\|_{L^p_t\left(I,L^q_x\right)} \le C \|\ga\|_{\FS^\al}.
\end{align}
\item\label{lemm3.2.2} For any $f \in L^{2p'}_t(I,L^{2q'}_x)$,
\begin{align}\label{HSS}
\left\|f A^\oc\right\|_{\FS^{2\al'}_{x \to (\tx)}} \le C'\|f\|_{L^{2p'}_t\left(I,L^{2q'}_x\right)}.
\end{align}
\item\label{lemm3.2.3} For any $f \in L^{2p'}_t(I,L^{2q'}_x)$,
\begin{align}\label{HSS-adjoint}
\left\|A^\os f\right\|_{\FS^{2\al'}_{(\tx) \to x}} \le C''\|f\|_{L^{2p'}_t\left(I,L^{2q'}_x\right)}.
\end{align}
Moreover, $\sqrt{C}$, $C'$ and $C''$ coincide.
\end{enumerate}
\end{lemm}

\begin{proof}
Assume~\eqref{lemm3.2.1}. Then we have
%transformation de plusieurs en un
\begin{equation}
\begin{split}
\left\|f A^\oc\right\|_\Sxtx^2 &= \left\|A^\os |f|^2 A^\oc\right\|_{\FS^{\al'}} \\
&= \sup \left\{ \left|\Tr\left[\ga A^\os |f|^2 A^\oc\right]\right | : \|\ga\|_{\FS^\al} \le 1\right\} \\
&= \sup \left \{ \abs{\int_I \Tr\left[A(t) \ga A(t)^* |f(t)|^2\right]dt } : \|\ga\|_{\FS^\al} \le 1 \right \} \\
&\le \sup \left \{ \|\rh(A(t)_\st \ga)\|_{L^p_t\left(I,L^q_x\right)} \|f\|_{L^{2p'}_t\left(I,L^{2q'}_x\right)}^2 : \|\ga\|_{\FS^\al} \le 1 \right \} \\
&\le C \|f\|_{L^{2p'}_t\left(I,L^{2q'}_x\right)}^2.
\end{split}
\end{equation}
Therefore, we obtain
\begin{align}
\left\|f A^\oc\right\|_\Sxtx \le \sqrt{C} \|f\|_{L^{2p'}_t\left(I,L^{2q'}_x\right)}.
\end{align}
We can prove that~\eqref{lemm3.2.2} yields $C \le (C')^2$ similarly. Since
\begin{align}
\left\|f A^\oc\right\|_{\Sxtx} = \left\|A^\os f\right\|_{\Stxx},
\end{align}
we completed the proof.
\end{proof}

The following lemma is sometimes useful:

\begin{lemm}\label{asym}
Let $I \subset \R$ be an interval. Assume that
\begin{align}
\|\rh(A_n(t)_\st \ga)\|_{L^{p_n}_t \left(I,L^{q_n}_x\right)} \le C_n \|\ga\|_{\FS^{\al_n}}
\end{align}
for $n=0,1$. Then we have
\begin{align}
\big\|\rh(A_0(t) \ga A_1(t)^*)\big\|_{L^p_t\left(I,L^q_x\right)} \le \sqrt{C_0 C_1} \|\ga\|_{\FS^\al},
\end{align}
where
\begin{align}
\frac{1}{X} = \frac{1}{2} \K{\frac{1}{X_0}+ \frac{1}{X_1}}, \quad X = p,q,\al.
\end{align}
\end{lemm}

\begin{proof}
Let $f \in L^{p'}_t(I,L^{q'}_x)$. $f$ can be decomposed as $f = f_0 f_1$, where
\begin{align}
\|f\|_{L^{p'}_t\left(I,L^{q'}_x\right)} = \|f_0\|_{L^{2p_0'}_t\left(I,L^{2q_0'}_x\right)} \|f_1\|_{L^{2p_1'}_t\left(I,L^{2q_1'}_x\right)}.
\end{align}
Then we have
\begin{multline}
\abs{\int_I \int_{\R^d} \rh(A_0(t)\ga A_1(t)^*)(x) f(t,x) dx dt}\\
\begin{aligned}
 &\leq \abs{\int_I \Tr (A_0(t) \ga A_1(t)^* f(t,x)) dt } \\
&\leq \abs{\int_I \Tr (A_1(t)^* f(t,x)A_0(t) \ga) dt} \\
&\leq \No{A_1^\os f A_0^\oc }_{\FS^{\al'}} \|\ga\|_{\FS^\al} \\
&\leq \|\ga\|_{\Sa} \left\|A_1^\os f_1\right\|_{\FS^{2\al_1'}_{(t,x) \to x}} \left\|f_0 A_0^\oc\right\|_{\FS^{2\al_0'}_{x\to(t,x)}}\\
&\leq \sqrt{C_0C_1}\|\ga\|_\Sa \|f\|_{L^{p'}_tL^{q'}_x}.\qedhere
\end{aligned}
\end{multline}
\end{proof}
\setcounter{theo}{5}
\setcounter{equation}{45}
\subsection{Propagators with time-dependent potentials} 

Let $p,q \in (1,\I)$ and $\al \in [1,\I]$. We call $(p,q,\al)$ admissible when $\frac{1}{\al} \ge \frac{1}{dp} + \frac{1}{q}$ and $\al < p$ hold. The following result is the most essential tool in this paper:

\begin{theo}[{\cite[Theorem~1]{2014FLLS}; \cite[Theorem~8]{2017FS}; \cite[Theorem~1.5]{2019BHLNS}}] \label{Ostri}
Let $d \ge 1$. Let $(p,q,\al)$ be admissible and $\si(p,q)=\frac{2}{p}+\frac{d}{q}=d-2s$ for $s \in (0,d/2)$. Then it holds that
\begin{align}
\|\rh(U(t)_\st Q_0)\|_{L^p_t L^q_x} \ls \|\sd^s_\st Q_0\|_{\FS^\al}.
\end{align}
\end{theo}


By Theorems~\ref{SCK}, \ref{Ostri} and Lemma~\ref{DP}, we obtain the following corollary.

\begin{coro}\label{DS}
Let $d \geq 1$. Let $(p_n,q_n,\al_n)$ be admissible and $\si(p_n,q_n)=d$ for $n = 0,1$. Define
\begin{align}
\CD(g_0,g_1) = \int_0^t g_0(t)U(t-\ta)g_1(\ta)d\ta.
\end{align}
Then we have
\begin{align}
\|\CD(g_0,g_1)\|_{\FS^{\al'}_{t,x}} \ls \prod_{n=0,1} \|g_n\|_{L^{2p_n'}_t L^{2q_n'}_x},
\end{align}
where $\frac{1}{\al'} = \frac{1}{2\al_0'} + \frac{1}{2\al_1'}$.
\end{coro}

\begin{proof}
First, we consider the operator
\begin{align}
\BT(g_0,g_1) = \int_0^\I g_0(t)U(t-\ta)g_1(\ta)d\ta = g_0 U^\oc U^\os g_1.
\end{align}
By Lemma~\ref{DP} and Theorem~\ref{Ostri}, we obtain
\begin{align}
\left\|\BT(g_0,g_1)\right\|_{\FS^{\al'}_{\tx}} &\le \left\|g_0 U^\oc\right\|_{\FS^{2\al_0'}_{x \to (\tx)}} \left\|U^\os g_1\right\|_{\FS^{2\al_1'}_{(\tx)\to x}} \\
&\ls \prod_{n=0,1} \|g_n\|_{L^{2p_n'}_t L^{2q_n'}_x}.
\end{align}
Therefore, we get the desired estimate by Theorem~\ref{SCK}.
\end{proof}
\section{Strichartz estimates for density functions}\label{sect4}

In this section, we prove the key estimates in this paper.
\begin{theo}\label{Key}
Let $d \ge 3$. Let $0 \le \wt{s} \le s \le 1$. Let $p,q,\mu,\nu,\al \in (1,\I)$ satisfy $\si(p,q)=d-s$ and $\si(\mu,\nu) = 2+\wt{s}$. Let $(p,q,\al)$ be admissible. Assume that there exist $p_j,q_j,\al_j \in (1,\I)$ for $j=0,1,2$ and $\mu_j,\nu_j \in (1,\I)$ for $j=0,1$ such that
\begin{align}
&\si(p_j,q_j) = d-j,\quad (p_j,q_j,\al_j) \text{ are admissible for } j=0,1,2, \\
&\si(\mu_0,\nu_0) = 2, \quad \si (\mu_1,\nu_1) = 2+s_1 := 2 + \frac{\wt{s}}{s}, \quad 1+s_1 < \frac{2}{\mu_1}+\frac{d}{2},\\
&\frac{1}{X} = \frac{1-s}{X_0} + \frac{s}{X_1}, \quad X = p,q,\mu,\nu,\al,
\quad \frac{1}{Y_1} = \frac{1}{2} \K{\frac{1}{Y_0} + \frac{1}{Y_2}}, \quad Y=p,q,\al.
\end{align}
\end{theo}
\setcounter{theo}{2}
\setcounter{equation}{10}
Before proving the above theorem, we define some notations. For $\vF=(F_1,\dots, F_n) \linebreak\in \CS(\BR^d)^n$, define
\begin{align}
\CW_n [\vF](t) = (-i)^n \int_0^t dt_1 \int_0^{t_1} dt_2 \cdots \int_0^{t_{n-1}} dt_n
\CU[F_1](t_1) \cdots \CU[F_n](t_n),
\end{align}
where $\CU[V](t):= U(t)^*V(t)U(t)$. And we define multilinear operators by
\begin{align}
T_{n,m}[\vF,Q_0, \vG](t) := \rh \big( U(t)\CW_n[F](t) Q_0 \CW_m [G](t)^* U(t)^* \big)
\end{align}
for $\vF \in \CS(\BR^d)^n$ and $\vG \in \CS(\BR^d)^m$. We have

\begin{lemm}\label{muliti est}
Under the same assumptions as in Theorem~\ref{Key},
The following multilinear estimates hold:
\begin{align}
&\|T_{n,m} [\vF,Q_0,\vG]\|_{L^p_t H^{s,q}_x}
\ls C_0^{n+m}
\prod_{j=1}^n \|F_j\|_{L^{\mu}_t H^{s,\nu}_x}
\|Q_0\|_{\CH^{s,\al}_x}
\prod_{k=1}^m \|G_k\|_{L^{\mu}_t H^{s,\nu}_x}.
\end{align}
\end{lemm}

\begin{proof}
It suffices to prove the following two estimates:
\begin{align}
\|T_{n,m} [\vF,Q_0,\vG]\|_{L^{p_0}_t L^{q_0}_x}
&\ls C_0^{n+m}
\prod_{j=1}^n \|F_j\|_{L^{\mu_0}_t L^{\nu_0}_x}
\|Q_0\|_{\FS^{\al_0}_x}
\prod_{k=1}^m \|G_k\|_{L^{\mu_0}_t L^{\nu_0}_x} \label{L2}
,
\\
\|T_{n,m} [\vF,Q_0,\vG]\|_{L^{p_1}_t H^{1,q_1}_x}
&\ls C_0^{n+m}
\prod_{j=1}^n \|F_j\|_{L^{\mu_1}_t H^{1,\nu_1}_x}
\|Q_0\|_{\CH^{1, \al_1}_x}
\prod_{k=1}^m \|G_k\|_{L^{\mu_1}_t H^{1,\nu_1}_x}. \label{H1}
\end{align}
We only prove~\eqref{H1} because~\eqref{L2} can be shown similarly. We have
\begin{multline}
\|T_{n,m}[\vF,Q_0,\vG]\|_{L^{p_1}_t H^{1,q_1}_x} \\
\sim \|T_{n,m}[\vF,Q_0,\vG]\|_{L^{p_1}_t L^{q_1}_x} + \|\na T_{n,m}[\vF,Q_0,\vG]\|_{L^{p_1}_t L^{q_1}_x}.
\end{multline}
Let $f \in L^{2p_1'}_t L^{2q_1'}_x$. Note that $\sd = \sdi - L \na$ with $L:= \na \sdi$. We have
%transfo de 2 pour un
\begin{equation}
\begin{split}
f(t)U(t)\CW_n[\vF](t) \sd^{-1} &= f(t)U(t)\sdi \left(\sdi - L\na\right)\CW_n[\vF](t) \sd^{-1} \\
&=:W_1 - W_2.
\end{split}
\end{equation}
We can assume $n \ge 1$ because when $n=0$, we obtain~\eqref{half} by Theorem~\ref{Ostri}. $W_1$ is a good term, and the problem is $W_2$. We have
\begin{multline}
W_2 = f(t)U(t)\sdi L \sum_{k=1}^n \CW_n[\vF^k](t) \sd^{-1} \\+ f(t)U(t)\sdi L \CW_n[\vF](t) \na \sd^{-1},
\end{multline}
where $\vF^k = (F_1, \cdots, F_{k-1}, \na F_k, F_{k+1}, \cdots, F_n)$. The last term is harmless. Let $F_{1,0}:= |F_1|^{1/2}$ and $F_1 = F_{1,0}F_{1,1}$. Note that $L^{\mu_1}_t H^{1,\nu_1}_x \hookrightarrow L^{\mu_1}_t L^{\wt{\nu_1}}_x$ with $\si(\mu_1, \wt{\nu_1}) =2$. Then we have
%transformations de plusieurs en un.
\begingroup
\allowdisplaybreaks
\begin{align}
&\Bigg\| f(t)U(t)\sdi L
\int_0^t dt_1
\int_0^{t_1} dt_2 \cdots \int_0^{t_{n-1}} dt_n
\CU[F_1](t) \\
&\mkern250mu\cdots \CU[\na F_k](t_k) \cdots \CU[F_n](t_n) \sdi\Bigg\|_{\FS^{2\al_1'}_{x \to (\tx)}} \nonumber\\
& \le \No{f(t)U(t)\sdi L \int_0^t dt_1 U(t_1)F_{1,0}(t_1)}_{\FS^{2\al_1'}_\tx} \nonumber\\
&\quad \times \bigg\| F_{1,1}(t_1)U(t_1)\int_0^{t_1} dt_2 \cdots \int_0^{t_{n-1}} dt_n
\CU[F_1](t) \cdots \CU[\na F_k](t_k)\nonumber\\
&\mkern250mu \cdots \CU[F_n](t_n) \sdi \bigg\|_{\CB\left(L^2_x \to L^2_\tx\right)} \nonumber\\
&\ls C_0
\|f\|_{L^{2p_1'}_t L^{2q_1'}_x}
\|F\|_{L^{\mu_1}_t L^{\wt{\nu_1}}_x}^{1/2} \bigg\| F_{1,1}(t_1)U(t_1)
\int_0^{t_1} dt_2 \CU[F_2](t_1)\nonumber\\
&\quad \times
\cdots
\int_0^{t_{k-1}} dt_k \CU[\na F_k](t_k)
\cdots
\int_0^{t_{n-1}} dt_n \CU[F_n](t_n) \sdi\bigg\|_{\CB\left(L^2_x \to L^2_\tx\right)} =: I.\nonumber
\end{align}
\endgroup
To justify the last inequality, we used Theorem~\ref{SCK}. Moreover, we have
%transfo de plusieurs pour un
\begin{multline}
I \le C_0^k
\|f\|_{L^{2p_1'}_t L^{2q_1'}_x}
\prod_{j=1}^{k-1}\|F_j\|_{L^{\mu_1}_t L^{\wt{\nu_1}}_x}\\
\begin{aligned}
&\times
\bigg\| (\na F_k)(t_k) U(t_k)
\int_0^{t_k}dt_{k+1} \CU[F_{k+1}](t_{k+1})
\cdots\\
&\mkern200mu\int_0^{t_{n-1}} dt_n \CU[F_n](t_n) \sdi \bigg\|_{\CB\left(L^2_x \to L^{\wt{r}'}_t L^{\wt{c}'}_x\right)} \\
\le&\, C_0^k \|f\|_{L^{2p_1'}_t L^{2q_1'}_x}
\prod_{j=1}^k \|F_j\|_{L^{\mu_1}_t H^{1,\nu_1}_x}\\
&\times
\No{ U(t_k)
\int_0^{t_k}dt_{k+1} \CU[F_{k+1}](t_{k+1})
\cdots
\int_0^{t_{n-1}} dt_n \CU[F_n](t_n) \sdi }_{\CB\left(L^2_x \to L^r_t L^{c^*}_x\right)} \\
\le&\, C_0^k \|f\|_{L^{2p_1'}_t L^{2q_1'}_x}
\prod_{j=1}^k \|F_j\|_{L^{\mu_1}_t H^{1,\nu_1}_x}\\
&\times
\No{
\sd U(t_k)
\int_0^{t_k}dt_{k+1} \CU[F_{k+1}](t_{k+1})
\cdots
\int_0^{t_{n-1}} dt_n \CU[F_n](t_n) \sdi }_{\CB\left(L^2_x \to L^r_t L^c_x\right)},
\end{aligned}
\end{multline}
where $\wt{r},\wt{c},r,c$ are defined as follows. Fix $r \ge 2$ such that $\tw-\frac{1}{\mu_1} < \frac{1}{r} < \min(\frac{d}{4}-\frac{s_1}{2}, 1-\frac{1}{\mu_1})$. Define $\wt{r} \in [2,\I]$ by $\frac{1}{\wt{r}'}=\frac{1}{\mu_1}+\frac{1}{r}$. Then there exists $c,\wt{c} \in [2,\I]$ such that $\si(r,c) = \si(\wt{r},\wt{c}) = \frac{d}{2}$. By direct calculation, we obtain
\begin{align}
\frac{1}{c^*}:=\frac{1}{\wt{c}'}-\frac{1}{\nu_1} = \frac{1}{c}-\frac{s_1}{d} > 0.
\end{align}
Since
\begin{align}
\No{
\int_0^t U(t-\ta)(g(\ta)h(\ta)) d\ta }_{L^r_t H^{1,c}_x}
\le C_0 \|g\|_{L^{\mu_1}_t H^{1,\nu_1}_x} \|h\|_{L^r_t H^{1,c}_x},
\end{align}
we have
%transfo de plusieurs en un
\begin{multline}
\No{
\sd U(t_k)
\int_0^{t_k}dt_{k+1} \CU[F_{k+1}](t_{k+1})
\cdots
\int_0^{t_{n-1}} dt_n \CU[F_n](t_n) \sdi }_{\CB\left(L^2_x \to L^r_t L^c_x\right)} \\
\le C_0^{n-k}
\prod_{j=k+1}^n \|F_j\|_{L^{\mu_1}_t H^{1,\nu_1}_x}.
\end{multline}
From the above, we obtain
\begin{align}
\left\|f(t) U(t) \sdi L \CW_n[\vF^k](t) \sd^{-1}\right\|_{\FS^{2\al_1'}_{x \to (\tx)}}
\ls C_0^n \|f\|_{L^{2p_1'}_t L^{2q_1'}_x}\prod_{j=1}^n \|F_j\|_{L^{\mu_1}_t H^{1,\nu_1}_x},
\end{align}
which yields
\begin{align}\label{half}
\left\|f(t)U(t)\CW_n[\vF](t) \sd^{-1} \right\|_{\FS^{2\al_1'}_{x \to (\tx)}}
\ls C_0^n \|f\|_{L^{2p_1'}_t L^{2q_1'}_x}\prod_{j=1}^n \|F_j\|_{L^{\mu_1}_t H^{1,\nu_1}_x}.
\end{align}
Therefore, we conclude that
\begin{multline}
\|T_{n,m}[\vF,Q_0,\vG]\|_{L^{p_1}_t L^{q_1}_x}
\\
\ls C_0^{n+m} \K{\prod_{j=1}^n
\|F_j\|_{L^{\mu_1}_t H^{1,\nu_1}_x}}
\|Q_0\|_{\CH^{1,\al_1}}
\K{\prod_{k=1}^m \|G_k\|_{L^{\mu_1}_t H^{1,\nu_1}_x}}.
\end{multline}
Finally, we estimate $\|\na T_{n,m}[\vF,Q_0,\vG]\|_{L^{p_1}_t L^{q_1}_x}$. In the completely same way as the above, we get
%transfo de 4 en 2
\begin{multline}\label{End H1}
\|T_{n,m} [\vF,Q_0,\vG]\|_{L^{p_2}_t L^{q_2}_x}\\
\ls C_0^{n+m}
\K{\prod_{j=1}^n \|F_j\|_{L^{\mu_1}_t H^{1,\nu_1}_x}}
\|Q_0\|_{\CH^{1,\al_2}}
\K{\prod_{k=1}^ m\|G_k\|_{L^{\mu_1}_t H^{1,\nu_1}_x}}, 
\end{multline}
\begin{multline}\label{deriv L2}
\left\|\rh\Dk{\na U(t)\CW_n[\vF](t) \sdi Q_0 \sdi \CW_n[\vG](t)^* U(t)^* \na} \right\|_{L^{p_0}_t L^{q_0}_x} \\
\ls C_0^{n+m} \K{\prod_{j=1}^n \|F_j\|_{L^{\mu_1}_t H^{1,\nu_1}_x}}
\|Q_0\|_{\FS^{\al_0}}
\K{\prod_{k=1}^m \|G_k\|_{L^{\mu_1}_t H^{1,\nu_1}_x}}. 
\end{multline}
Therefore, we obtain by Lemma~\ref{asym}
%transfo de 4 en 1
\begin{multline}
\left\|\na T_{n,m}[\vF,Q_0,\vG]\right\|_{L^{p_1}_t L^{q_1}_x} \\
\begin{aligned}
&\le \No{\rh \Dk{\na U(t) \CW_n[\vF](t) Q_0 \CW_m[\vG](t)^* U(t)}}_{L^{p_1}_t L^{q_1}_x} \\
&+ \No{\rh \Dk{U(t) \CW_n[\vF](t) Q_0 \CW_m[\vG](t)^* U(t) \na}}_{L^{p_1}_t L^{q_1}_x} \\
&\ls C_0^{n+m} \K{\prod_{j=1}^n \|F_j\|_{L^{\mu_1}_t H^{1,\nu_1}_x}} \|Q_0\|_{\CH^{1,\al_1}} \K{\prod_{k=1}^m \|G_k\|_{L^{\mu_1}_t H^{1,\nu_1}_x}}.\qedhere
\end{aligned}
\end{multline}
\end{proof}

\begin{thebibliography}{FLLS14}
\bibitem[BHL$^{+}$19]{2019BHLNS} Bez, Neal; Hong, Younghun; Lee, Sanghyuk; Nakamura, Shohei; Sawano, Yoshihiro On the Strichartz estimates for orthonormal systems of initial data with regularity, Adv. Math., Volume 354 (2019), 106736, 37 pages.
\bibitem[BS03]{2003BS} Birman, Mikhail Sh.; Solomyak, Michael Double operator integrals in a Hilbert space. Integral Equations Operator Theory, Integral Equations Oper. Theory, Volume 47 (2003) no. 2, pp. 131-168.
\bibitem[CK01]{2001CK} Christ, Michael; Kiselev, Alexander Maximal functions associated to filtrations, J. Funct. Anal., Volume 179 (2001) no. 2, pp. 409-425.
\bibitem[FLLS14]{2014FLLS} Frank, Rupert L.; Lewin, Mathieu; Lieb, Elliott H.; Seiringer, Robert Strichartz inequality for orthonormal functions, J. Eur. Math. Soc., Volume 16 (2014) no. 7, pp. 1507-1526.
\bibitem[FS17]{2017FS} Frank, Rupert L.; Sabin, Julien Restriction theorems for orthonormal functions, Strichartz inequalities, and uniform Sobolev estimates, Am. J. Math., Volume 139 (2017) no. 6, pp. 1649-1691.
\bibitem[GK70]{1970Go-Kre} Gohberg, Israel C.; Krein, Mark G. Theory and applications of Volterra operators in Hilbert space, Translations of Mathematical Monographs, 24, American Mathematical Society, 1970.
\bibitem[Sim05]{2005Simon} Barry Simon, Trace ideals and their applications, second edition, Mathematical Surveys and Monographs, vol.120, American Mathematical Society,2005.
\end{thebibliography}
\end{document}
